\section{问题设定与记号}

设模型参数由矩阵块组成：
\[
\theta=(W_1,\ldots,W_L),\qquad W_\ell\in\R^{m_\ell\times n_\ell}.
\]
对一个 batch，模型输出为
\[
F(\theta)\in\R^{c\times B},
\]
损失为
\[
\Loss(\theta)=\phi(F(\theta)).
\]
令
\[
G_\ell=\nabla_{W_\ell}\Loss(\theta)
\]
为第 \(\ell\) 个矩阵块的梯度。

若某层输入 activation matrix 为
\[
A_{\ell-1}=[h_{\ell-1}(x_1),\ldots,h_{\ell-1}(x_B)]\in\R^{n_\ell\times B},
\]
则该层 pre-activation 为
\[
Z_\ell=W_\ell A_{\ell-1}.
\]
设正下降更新为
\[
D_\ell=W_\ell-W_\ell^+.
\]
那么固定输入激活时的直接 pre-activation 变化为
\[
\Delta Z_{\ell,\mathrm{direct}}=-D_\ell A_{\ell-1}.
\]
如果激活函数 \(\sigma\) 是 \(L_\sigma\)-Lipschitz，则
\[
\|\Delta H_\ell\|_F
\le
L_\sigma\|D_\ell A_{\ell-1}\|_F.
\]
因此，本文将 \(D_\ell A_{\ell-1}\) 视为最基本的 activation perturbation object。

对矩阵 \(M\)，记
\[
\|M\|_{\op}=\sigma_1(M),\qquad
\|M\|_F=\left(\sum_i\sigma_i(M)^2\right)^{1/2},
\qquad
\|M\|_*=\sum_i\sigma_i(M).
\]
稳定秩定义为
\[
\srank(M)=\frac{\|M\|_F^2}{\|M\|_{\op}^2},
\]
梯度核秩定义为
\[
\nrank(G)=\frac{\|G\|_*^2}{\|G\|_F^2}.
\]
